\begin{circuitikz}[
    american, 
    >={Latex[length=5mm, width=2mm]},
    alt={A technical diagram illustrating a hardness test or indentation process. A wedge-shaped indenter with a trapezoidal body, subjected to an upward force vector $F$, is aligned above two hatched rectangular material layers. In the lower layer, a diamond-shaped indentation is marked with orthogonal dimensions $d_1$ and $d_2$, and dashed red lines project the corners of the upper indenter down to the indentation footprint.}
    ]

    \def\plateHeight{2cm}
    \def\plateWidth{6cm}
    \def\plateDepth{4cm}
    \def\indenterWidth{3cm}
    \def\bodyHeight{1cm}
    %\pgfmathsetlengthmacro{\longSide}{\indenterWidth + 2 * \indenterWidth * tan(30)}
    \pgfmathsetlengthmacro{\bodyWidth}{\indenterWidth + 2 * \bodyHeight * tan(60)}
    \pgfmathsetlengthmacro{\dentDiag}{\indenterWidth * cos(45)}
    
    \begin{pgfonlayer}{ft}

        \node[
            draw=BakerBlack,
            fill=HoneyYellow,
            isosceles triangle,
            isosceles triangle apex angle=136,
            minimum width=\indenterWidth,
            inner sep=0,
            shape border rotate=-90,
            anchor=lower side
        ](diamond)
        at (0,0)
        {};

        \draw[
            BakerBlack
        ]
        (diamond.apex)
        -- (diamond.lower side);

        %\tikzset{trapezium stretches=true}
        \node[
            draw=BakerBlack,
            fill=HoneyYellow,
            trapezium,
            minimum width=\bodyWidth,
            %minimum height=\bodyHeight,
            %inner sep=0,
            shape border rotate=180,
            anchor=top side
        ](body)
        at (diamond.lower side)
        {};

    \end{pgfonlayer}

    \begin{pgfonlayer}{bg}
        
        \node[
            draw=BakerBlack,
            preaction={
                fill=Iron1,
                %fill opacity=0.5,
            },
            pattern={Lines[
                angle=45,
                distance=0.25cm
            ]},
            pattern color=Abbey,
            minimum height=\plateHeight,
            minimum width=\plateWidth,
            anchor=north
        ](plate1)
        at (diamond.lower side)
        {};

        \node[
            draw=BakerBlack,
            preaction={
                fill=Iron1,
                %fill opacity=0.5,
            },
            pattern={Lines[
                angle=45,
                distance=0.25cm
            ]},
            pattern color=Abbey,
            minimum height=\plateDepth,
            minimum width=\plateWidth,
            anchor=north
        ](plate2)
        at ([yshift=-1cm]plate1.south)
        {};

    \end{pgfonlayer}   

    \begin{pgfonlayer}{main}

        \node[
            draw=BakerBlack,
            fill=Iron1,
            diamond,
            minimum height=\indenterWidth,            
            minimum width=\indenterWidth,
            inner sep=0,
        ](dent)
        at (plate2.center)
        {};

        \draw[
            BakerADARed,
            dashed
        ]
        (diamond.left corner)
        -- (dent.east);

        \draw[
            BakerADARed,
            dashed
        ]
        (diamond.right corner)
        -- (dent.west);

        \draw[
            BakerBlack,
            <->
        ]
        (dent.east)
        -- (dent.west)
        node[
            left,
            text=BakerBlack,
            fill=Iron1
        ]{$d_1$};

        \draw[
            BakerBlack,
            <->
        ]
        (dent.north)
        -- (dent.south)
        node[
            below,
            text=BakerBlack,
            fill=Iron1
        ]{$d_2$};

        \draw[
            UpNorthGreen,
            line width=3pt,
            >={Latex[length=7mm, width=4mm]},
            <-
        ]
        (body.bottom side)
        -- ++(0,1.5)
        node[
            text=BakerBlack,
            anchor=south
        ]{\Large{$F$}};

    \end{pgfonlayer}

    

     

\end{circuitikz}